\scanpage{67}
\deleted{\unsure{Comme|Corol}}
\deleted{On sait que}
\deleted{On sait que}

On a d'autre part, si
\[
\deleted{\unsure{0|\cdot}\unsure{\longrightarrow|\mathbin{-}}\Gamma\longrightarrow}\qquad \unsure{j_0|j_a}:\Gamma\longrightarrow\operatorname{Pic}(X_0)\qquad\deleted{\unsure{est|\longrightarrow}}
\]
est le composé, d'où des $j_n:\Gamma_n\mathbin{-}\operatorname{Pic}(X_0)_n$, \deleted{\unsure{donc} \unsure{l'inclusion}} \inserted{\unsure{on a}} l'inclusion \unsure{évidente|evidente}\unsure{:|;}
\[
\operatorname{Ker}i_n\subset\operatorname{Ker}j_n
\]
\margin{\uline{Lemme.} $\operatorname{Pic}(X)_n\hookrightarrow\operatorname{Pic}(X_0)_n$ (du moins si $S$ complet, ce qui suffit pour le th.) faisable par la rigidité}
Cette inclusion est \inserted{en fait} une \uline{identité}. Pour le voir, on note que si $\underline L$ est un faisceau inv. sur $X$, et si $\underline L_n=\underline L|X_n$ est \unsure{écrit|écris} comme puissance $n$ième, alors l'obstruction à écrire $\underline L_{n\unsure{+}1}$ comme puissance $n$ième est nulle (cf \unsure{FGA|PCA}, exposé sur Picard). On a d'autre part
\[
\operatorname{Pic}(X_0)=\underline{\operatorname{Pic}}_{X_0/k}(k)
\]
grâce au fait que $X_0$ a un $0$-cycle de degré $1$, d'où
\diagram{\begin{tikzcd}[column sep=large]
0\arrow[r,dash]&\operatorname{Pic}^{0}(X_0)\arrow[r,dash]\arrow[d,equal,"\mathrm{dfn}"']&\operatorname{Pic}(X_0)\arrow[r]&\operatorname{NS}(X_0)'\arrow[r]\arrow[d,equal,"\mathrm{dfn}"']&0\\
&\underline{\operatorname{Pic}}^{0}_{X_0/k}(k)&&\operatorname{Im}\bigl(\underline{\operatorname{Pic}}_{X_0/k}(k)\longrightarrow\underline{\operatorname{NS}}_{X_0/k}(k)\bigr)
\end{tikzcd}}
et $\operatorname{Pic}^{0}(X_0)$ est $n$-divisible, donc on a
\[
\operatorname{Pic}(X_0)_n\hookrightarrow\operatorname{NS}(X_0)_n
\]
\margin{\uline{joindre au lemme précédent}}
donc si
\[
\unsure{k}:\Gamma\longrightarrow\operatorname{NS}(X_0)
\]
est le composé $\Gamma\xrightarrow{j}\operatorname{Pic}(X_0)\longrightarrow\operatorname{NS}(X_0)$, on trouve
\[
\operatorname{Coker}u_n\simeq\operatorname{Ker}\unsure{k_n}
\]
Je dis que \inserted{$\Gamma$ \uline{est sans torsion et que}} $\delta$ est \uline{injectif} i.e. qu'un faisceau inv. sur $X$ qui est banal sur $X_\eta$, et qui sur $X_0$ est \inserted{num. équiv.} \deleted{\inserted{(comme \illegible{1} \unsure{de} torsion)}} à zéro, est trivial.
\margin{cf plus bas une démonstration}
\deleted{L'hyp. que l'image}
\deleted{$\unsure{L'\in\operatorname{Pic}(X)}$ \unsure{dans} $\operatorname{NS}(X_0)$ est nulle \deleted{\unsure{implique}} \inserted{\unsure{implique}} \deleted{\unsure{par le}} \unsure{indique|implique} (par rigidité) que que $\underline L$ est \unsure{inft|infiniment} $\ell$-divisible pour tt $\ell\ne p$, et on peut faire la division de façon \underline{unique}, en restant dans}
\unsure{finalt|finalement} il est banal, on trouve
une suite exacte
\deleted{\unsure{k|\xi}}
\deleted{\diagram{\begin{tikzcd}
\unsure{0|\cdot}\arrow[r,dash]&\Gamma\arrow[r,"\unsure{k}",dash]&\operatorname{NS}(X_0)'\arrow[r]&\unsure{\delta|\operatorname{Coker}k}\arrow[r]\arrow[d,equal]&0\\
&&&\operatorname{Coker}\unsure{k}\unsure{=\operatorname{Coker}|=}\unsure{\Delta|\delta}
\end{tikzcd}}}
\deleted{\inserted{$\operatorname{Coker}k=\operatorname{Coker}\Delta$}}
\margin{\inserted{\diagram{\begin{tikzcd}[column sep=small]
0\arrow[r]&\Gamma\arrow[r,"k'"]&\operatorname{NS}(X_0)'\arrow[r,"\lambda"]\arrow[d,equal]&\delta\arrow[r]&0\\
&&\text{\unsure{ss-grpe|ss-groupe} \unsure{des|de} }\operatorname{NS}(X_0)\text{ dont un multiple tombe ds }\Gamma
\end{tikzcd}}}}
d'où \inserted{par} \unsure{la|le} suite exacte habituelle
\[
\deleted{\unsure{\operatorname{Coker}|donc}}\quad \operatorname{Ker}k_n=\operatorname{Ker}k'_n\simeq\operatorname{Coker}{}_n\lambda
\]
Comme $\Gamma$ est sans torsion, \deleted{\unsure{donc|on}} $\delta'\xrightarrow{\lambda}\delta$ est injectif, et on trouve que $\operatorname{Coker}{}_n\lambda={}_n\delta/{}_n\delta'$. On termine grâce à la suite
\margin{66}
